%% SCRAP: scratch/COMMUNICATION_GUIDE
%% SOURCE: docs/working/scratch/COMMUNICATION_GUIDE.adoc
%% STATUS: HISTORICAL
%% FITS: none
%% EDITORIAL: lifted — prose rewritten to press voice

\section{StarshipOS Analogy Reference}

StarshipOS's physics-driven adaptive runtime admits a layered set of analogies that
allow the same technical concept to be communicated across audiences from primary school
through academic research. The anchoring fact in every case is that the system implements
verifiable feedback loops: guard proofs, telemetry fields, and runtime invariants are
concrete, not metaphorical.

\begin{itemize}
  \item \textbf{Elementary.} The OS watches when activity grows too noisy or energy runs
        low and nudges the system back toward calm, the way a classroom helper restores
        order.
  \item \textbf{Middle school.} Health and energy bars on a video game HUD track
        \emph{heat} and \emph{noise}; the OS reads those bars and adjusts tactics to
        keep them balanced.
  \item \textbf{High school.} A car dashboard combines speed, temperature, and fuel
        readings; cruise control adjusts throttle and cooling to smooth the ride. The OS
        applies the same pattern to compute metrics.
  \item \textbf{Trade professional (HVAC / facilities).} A building management controller
        reads pressure and temperature, then modulates fans and valves for steady airflow.
        The OS treats execution metrics identically.
  \item \textbf{University STEM.} Voltage, current, and heat readings drive a lab
        feedback controller that keeps experiments within bounds. StarshipOS treats VM
        metrics as the equivalent sensor inputs.
  \item \textbf{Working engineer / developer.} The runtime implements a PID loop around
        the VM, exposing entropy, latency, and heat as signals, enforcing invariants via
        machine-checked proofs, and delegating tuning to policy modules.
  \item \textbf{Researcher / academic.} The system models state variables as
        thermodynamic constructs, verifies semantics in Isabelle/HOL, and runs scheduling
        and policy on formally constrained, measurable signals within a dynamical-systems
        framework.
\end{itemize}

The goal is to move effortlessly between these framings depending on the audience while
keeping every analogy grounded in the same set of concrete artefacts.
